\documentclass[11pt,a4paper]{article}
\usepackage[margin=0.85in,top=1in,bottom=1in]{geometry}
\usepackage{amsmath,amssymb,amsfonts}
\usepackage{graphicx}
\usepackage{float}
\usepackage{enumitem}
\usepackage{microtype}
\usepackage{caption}
\usepackage[hidelinks]{hyperref}
\usepackage[most,breakable]{tcolorbox}
\tcbuselibrary{breakable,skins}
\usepackage{fontspec}
\setmainfont{DejaVu Serif}
\setsansfont{DejaVu Sans}
\setmonofont{DejaVu Sans Mono}
\captionsetup{font=small,labelfont=bf,skip=4pt}
\setlist{leftmargin=1.5em,topsep=2pt}

%% Breakable problem card — prevents content cutoff across pages
\newtcolorbox{solcard}[3]{
  breakable, enhanced,
  colback=white, colframe=gray!50, arc=2pt,
  title={\textbf{#1}\hfill\textsf{\small #3\,pts}},
  fonttitle=\normalsize\sffamily\bfseries,
  before upper={\small\textit{#2}\par\smallskip\hrule height 0.4pt\smallskip},
  top=4pt, bottom=6pt, left=7pt, right=7pt,
  parbox=false,
}

\title{\textbf{Quartz resonator --- Solution}}
\author{\normalsize X26 (2026) — English translation}
\date{}
\begin{document}\maketitle\vspace{-0.5em}
\begin{solcard}{A0}{Write down the code for the package containing the electronic components on your answer sheets. Check that it is complete.}{??}

To complete this step, you need to carefully look around the entire workplace. If you look carefully enough, there is a high probability of finding a package with electronic components on the table, on which an individual code in the “$OOO~00~00$” format is applied. After discovering the package, we carefully rewrite the entire code using a pen into a specially designated place on the answer sheet (this place is easily identified by the inscription “A0” on the left). At this stage you need to be extremely focused, otherwise it may lead to mistakes. The most common mistakes are: rewriting only the last 4 characters (“$00~00$”); rewriting characters in the wrong order (“$00~00~OOO$”); errors in spelling of characters (“$OOO~oo~oo$”); adding extra characters (“$00~00~(OOO)$”); no code on the answer sheet. The next step in completing this item is to check the completeness of the package. The check procedure is as follows: Open the bag. Carefully pour out all the ingredients from the bag. Make sure the bag is really empty. Familiarize yourself with the list of equipment from the condition (if you are sufficiently insinuating when reading the condition of the problem, it can be found by the inscription Equipment). Check for the presence of a quartz resonator in the poured elements (looks like a gray box with legs). Read the table of capacitor markings presented after the photo of the equipment in the condition. Attention! It is required to compare in order the presence of capacitors in the table (see point 3) and on the table, poured out of the bag in point 1 (see point 1). If the mapping between the names in the table and the capacitors on the table is one-to-one, then you can proceed to the next step. Otherwise, you should contact the duty officer.


To complete this step, you need to carefully look around the entire workplace. If you look carefully enough, there is a high probability of finding a package with electronic components on the table, on which an individual code in the “$OOO~00~00$” format is applied.


After discovering the package, we carefully rewrite the entire code using a pen into a specially designated place on the answer sheet (this place is easily identified by the inscription “A0” on the left). At this stage you need to be extremely focused, otherwise it may lead to mistakes. The most common mistakes are:


The next step in completing this item is to check the completeness of the package. The verification procedure is as follows:

\par\smallskip\noindent\textbf{Answer:} \textbackslash{}textbf{Answer:} \[OOO~00~00\]
\end{solcard}

\begin{solcard}{A1}{Find the signal frequency $f_{\max}$ at which the voltage amplitude across the quartz resonator is maximum. The difference between $f_{\max}$ and $2~\mathrm{MHz}$ is no more than $1~\mathrm{kHz}$.}{0.20}
\par\smallskip\noindent\textbf{Answer:} \textbackslash{}textbf{Answer:} \[f_{\max} = 1.999860 ~\mathrm{MHz}.\]
\end{solcard}

\begin{solcard}{A2}{Record the ratio of voltage amplitudes on two oscilloscope channels $\dfrac{U_{\text{CH}2}}{U_{\text{CH}1}}$ at signal frequency $f_{\max}$.}{0.05}
\par\smallskip\noindent\textbf{Answer:} \textbackslash{}textbf{Answer:} \[\dfrac{U_{\text{CH}2}}{U_{\text{CH}1}}=0.733.\]
\end{solcard}

\begin{solcard}{A3}{Remove the dependence of $2U_{\text{CH}2}$ and $\varphi$ on the frequency $f$ in the range $[f_{\max}-250~\mathrm{Hz}; f_{\max}+250~\mathrm{Hz}]$, where $\varphi$ is the phase difference between the signals from channels 2 and 1. Measure at least 30 points from the given frequency range.}{2.80}
\par\smallskip\noindent\textbf{Answer:} \textbackslash{}textbf{Answer:} $f, ~\mathrm{MHz}$ $U_{\text{CH}2},~ мВ$ $\Delta t, ~\mathrm{ns}$ $\varphi, ~ рад$ $|Z|, ~\mathrm{\Omega}$ $f, ~\mathrm{MHz}$ $U_{\text{CH}2},~ мВ$ $\Delta t, ~\mathrm{ns}$ $\varphi, ~ рад$ $|Z|, ~\mathrm{\Omega}$ 1.999610 80.0 72 0.905 852 1.999850 372 0 0.000 12482 1.999630 77.6 72 0.905 825 1.999855 376 12 0.151 11808 1.999650 72.0 68 0.854 765 1.999858 376 16 0.201 11101 1.999660 64.8 88 1.106 663 1.999860 384 26 0.327 9685 1.999670 58.8 84 1.055 603 1.999865 380 32 0.402 8458 1.999680 55.6 84 1.055 568 1.999870 376 42 0.528 6978 1.999690 52.0 84 1.055 530 1.999875 364 44 0.553 6466 1.999700 49.2 80 1.005 502 1.999880 352 58 0.729 5113 1.999710 44.8 78 0.980 457 1.999890 332 64 0.804 4471 1.999720 41.2 64 0.804 423 1.999900 312 74 0.930 3788 1.999730 36.4 56 0.704 374 1.999910 292 78 0.980 3422 1.999740 34.8 34 0.427 360 1.999920 276 82 1.030 3132 1.999750 35.6 0 0.000 372 1.999930 258 80 1.005 2970 1.999760 39.6 -26 -0.327 415 1.999940 246 84 1.056 2751 1.999770 50.4 -38 -0.477 536 1.999950 236 84 1.056 2637 1.999780 67.2 -50 -0.628 727 1.999960 228 86 1.081 2513 1.999790 88.8 -58 -0.729 981 1.999970 220 86 1.081 2423 1.999800 134 -54 -0.679 1597 1.999980 212 90 1.131 2280 1.999810 158 -56 -0.704 1936 1.999990 208 86 1.081 2287 1.999815 170 -52 -0.653 2154 2.000000 202 88 1.106 2194 1.999820 206 -54 -0.679 2722 2.000010 200 88 1.106 2172 1.999825 222 -44 -0.553 3191 2.000020 196 92 1.156 2083 1.999830 266 -49 -0.616 3953 2.000030 192 92 1.156 2040 1.999835 284 -32 -0.402 5079 2.000040 190 92 1.156 2019 1.999838 302 -22 -0.276 6221 2.000050 188 92 1.156 1997 1.999840 322 -20 -0.251 7243 2.000070 182 80 1.005 2048 1.999843 332 -14 -0.176 8239 2.000090 176 76 0.955 2012 1.999845 342 -10 -0.126 9199 2.000110 174 80 1.005 1950 1.999848 356 -6 -0.075 10610
\end{solcard}

\begin{solcard}{A4}{Plot a graph of the impedance modulus of the quartz resonator $|Z|$ versus frequency $f$ over the measured range.}{2.50}

Let's introduce notations for the voltages that we measure on an oscilloscope: \[U_0 \equiv U_{\text{CH}1};\quad U_1\equiv U_{\text{CH}2}.\] Let's write the equations for amplitudes and phases for two channels of the oscilloscope: \[\tilde U_{\text{CH}1}= \tilde I(R_0 + Z);\] \[\tilde  U_{\text{CH}2} = \tilde I Z.\] Divide the equations and get: \[\dfrac{U_{\text{CH}2}}{U_{\text{CH}1}} e^{i\varphi}= \dfrac{Z}{R_0+Z}.\] Now we can get an expression for the impedance of a quartz resonator: \[Z = U_{\text{CH}2}R_0\cdot\dfrac{e^{i\varphi}}{U_{\text{CH}1}-U_{\text{CH}2}e^{i\varphi}}\] This means that the impedance module can be recalculated using the following formula: \[|Z| = \dfrac{U_{\text{CH}2}R_0}{\sqrt{(U_{\text{CH}1}-U_{\text{CH}2}\cos\varphi)^2+(U_{\text{CH}2}\sin\varphi)^2}}\] Let's recalculate (see solution A3) and plot the dependence. \[f_0=1.999000~\mathrm{MHz}\]


\[\tilde U_{\text{CH}1}= \tilde I(R_0 + Z);\]


\[\tilde  U_{\text{CH}2} = \tilde I Z.\] Let's divide the equations and get: \[\dfrac{U_{\text{CH}2}}{U_{\text{CH}1}} e^{i\varphi}= \dfrac{Z}{R_0+Z}.\] Now we can get an expression for the impedance of a quartz resonator: \[Z = U_{\text{CH}2}R_0\cdot\dfrac{e^{i\varphi}}{U_{\text{CH}1}-U_{\text{CH}2}e^{i\varphi}}\] This means that the impedance module can be recalculated using the following formula: \[|Z| = \dfrac{U_{\text{CH}2}R_0}{\sqrt{(U_{\text{CH}1}-U_{\text{CH}2}\cos\varphi)^2+(U_{\text{CH}2}\sin\varphi)^2}}\] Let's perform the recalculation (see solution A3) and plot the dependence.


\[f_0=1.999000~\mathrm{MHz}\]

\begin{figure}[H]
  \centering
  \includegraphics[width=0.72\textwidth]{E3_S_EN_assets/fig1.jpg}
\end{figure}
\end{solcard}

\begin{solcard}{A5}{From the plotted graph $|Z|(f)$, find the frequencies $f_{рез}$ and $f_{анти}$ at which $|Z|$ is maximum and minimum, respectively.}{0.10}
\par\smallskip\noindent\textbf{Answer:} \textbackslash{}textbf{Answer:} \[f_{рез} = 1.999850~\mathrm{MHz}; \quad f_{анти} = 1.999740~\mathrm{MHz}.\]
\end{solcard}

\begin{solcard}{A6}{Express $\omega_s, \omega_p, \alpha$ in terms of $R, L, C, C_0$.}{0.30}

Impedance of parallel connection of capacitor $C_0$ and $RLC$-chain: \[\dfrac{1}{Z} = \dfrac{1}{R+i\omega L + \frac{1}{i\omega C}} + i\omega C_0 \] Let's simplify the expression: \[\dfrac{1}{Z} = \dfrac{1+i\omega C_0 \left(R+i\omega L + \frac{1}{i\omega C}\right)}{R+i\omega L + \frac{1}{i\omega C}}\]\[Z(\omega) = \dfrac{\frac{1}{LC}-\omega^2 + i\omega \frac{R}{L}}{i\omega C_0 \left(\frac{1}{LC_0}+\frac{1}{LC}-\omega^2 +i\omega\frac{R}{L}\right)}\]

\par\smallskip\noindent\textbf{Answer:} \textbackslash{}textbf{Answer:} Thus: \[\omega_s = \sqrt{\dfrac{1}{LC}};\quad \omega_p = \sqrt{\dfrac{1}{LC}\left(1+\dfrac{C}{C_0}\right)}; \quad \alpha = \dfrac{R}{L}.\]
\end{solcard}

\begin{solcard}{A7}{Using the approximations described above, express $|Z|^2(\omega)$ in terms of $\alpha, \omega, \omega_s, \omega_p, C_0$.}{0.10}

Let us take advantage of the fact that for any two complex numbers $z_1, z_2$ $\left| \frac{z_1}{z_2}\right| = \frac{|z_1|}{|z_2|}$ holds:

\par\smallskip\noindent\textbf{Answer:} \textbackslash{}textbf{Answer:} \[ |Z|^2 = \dfrac{(\omega_s^2 - \omega^2)^2 + (\omega_p \alpha)^2}{(\omega_p C_0)^2 ((\omega_p^2 - \omega^2)^2 + (\omega_p \alpha)^2)} .\]
\end{solcard}

\begin{solcard}{A8}{Using the approximations described above, express $f_{рез}$ and $f_{анти}$ in terms of $\omega_s, \omega_p$.}{0.20}

\[\dfrac{\mathrm{d}|Z|^2}{\mathrm{d}\omega^2} = \dfrac{2}{(\omega_p C_0)^2}\cdot\dfrac{-(\omega_s^2-\omega^2)((\omega_p^2 - \omega^2)^2 + (\omega_p \alpha)^2) + (\omega_p^2-\omega^2)((\omega_s^2 - \omega^2)^2 + (\omega_p \alpha)^2)}{((\omega_p^2 - \omega^2)^2 + (\omega_p \alpha)^2)^2}\]\[\dfrac{\mathrm{d}|Z|^2}{\mathrm{d}\omega^2} = \dfrac{2}{(\omega_p C_0)^2 } \dfrac{(\omega_p^2 - \omega_s^2)((\omega_p \alpha)^2-(\omega_p^2 - \omega^2)(\omega_s^2 - \omega^2)) }{((\omega_p^2 - \omega^2)^2 + (\omega_p \alpha)^2)^2}\]\[\dfrac{\mathrm{d}|Z|}{\mathrm{d}\omega} = \dfrac{2|Z|}{\omega(\omega_p C_0)^2 } \dfrac{(\omega_p^2 - \omega_s^2)((\omega_p \alpha)^2-(\omega_p^2 - \omega^2)(\omega_s^2 - \omega^2)) }{((\omega_p^2 - \omega^2)^2 + (\omega_p \alpha)^2)^2}\]Let us find the frequencies at which the derivative $|Z|$ is reset to zero: \[(\omega_p \alpha)^2-(\omega_p^2 - \omega^2)(\omega_s^2 - \omega^2) = 0\]\[-\omega^4 + \omega^2 (\omega_p^2+\omega_s^2) + (\omega_p\alpha )^2-(\omega_p\omega_s)^2=0.\]Solution of the quadratic equation: \[\omega = \sqrt{\dfrac{(\omega_s^2+\omega_p^2)\pm \sqrt{(\omega_s^2+\omega_p^2)^2+4\omega_p^2 (\alpha^2 - \omega_s^2)} }{2}}\]Using the approximation $\alpha\ll\omega_s$, we obtain the frequencies at which the derivative $|Z|$ is reset: $\omega = \omega_s$ and $\omega =\omega_p$.


Note. In paragraph A12 we will show the validity of this neglect for a given quartz resonator.


Now let’s determine which of the two points has resonance and which has antiresonance: \[|Z|^2 (\omega_p) = \dfrac{1}{(\omega_p C_0)^2}\cdot\dfrac{(\omega_s^2 - \omega_p^2)^2 + (\omega_p \alpha)^2}{ (\omega_p \alpha)^2}\]\[|Z|^2 (\omega_s) = \dfrac{1}{(\omega_p C_0)^2}\cdot\dfrac{ (\omega_p \alpha)^2}{(\omega_s^2 - \omega_p^2)^2 + (\omega_p \alpha)^2}\]Now we can see that at $\omega = \omega_s$ there will be antiresonance, and at $\omega = \omega_p$ there will be resonance.

\par\smallskip\noindent\textbf{Answer:} \textbackslash{}textbf{Answer:} \[f_{рез} = \dfrac{\omega_p}{2\pi}; \quad f_{анти} = \dfrac{\omega_s}{2\pi}\]
\end{solcard}

\begin{solcard}{A9}{Remove the dependence of the resonance frequencies $f_{рез}$ and antiresonance $f_{анти}$ of the quartz resonator on the external capacitance $C_{\text{ext}}$. As an external capacitance $C_{\text{ext}}$, it is enough to use each of the provided capacitors separately, without combining them.}{0.45}
\par\smallskip\noindent\textbf{Answer:} \textbackslash{}textbf{Answer:} $C_{\text{ext}}$, pF $f_{рез}$, MHz $f_{анти}$, MHz 0 1.999850 1.999750 6.0 1.999842 1.999755 15 1.999832 1.999755 10 1.999837 1.999755 27 1.999822 1.999755 33 1.999816 1.999755 47 1.999807 1.999755 56 1.999804 1.999755 68 1.999798 1.999755 82 1.999794 1.999755
\end{solcard}

\begin{solcard}{A10}{Construct a linearized graph of $C_{\text{ext}}$ versus $f_{рез}$ and $f_{анти}$.}{0.80}

Let us write the expression for the resonant frequency of a quartz resonator with a parallel connection $C_{\text{ext}}$: \[\omega_p(C_{\text{ext}}) = \sqrt{\dfrac{1}{LC}\left(1+\dfrac{C}{C_0+C_{\text{ext}}}\right)} = \omega_s\sqrt{1+\dfrac{C}{C_0+C_{\text{ext}}}}\]Then after the transformation we can obtain linearization: \[C_{{ext}} = \dfrac{f_{анти}^2 C}{f_{рез}^2 - f_{анти}^2} - C_0\]

\begin{figure}[H]
  \centering
  \includegraphics[width=0.72\textwidth]{E3_S_EN_assets/fig2.jpg}
\end{figure}
\end{solcard}

\begin{solcard}{A11}{From the results of points A6, A10, obtain the values ​​$L,C,C_0$.}{1.30}

From the linearized graph in A10 we obtain the slope and intercept, which are equal to: \[C = 5.3 ~фФ;\] \[C_0 = 53.9~\mathrm{pF}.\] Now, knowing the value of $f_{анти} = 1.999740~\mathrm{MHz}$, we can obtain the value of $L$: \[L = \dfrac{1}{4 \pi^2 f_\text{анти}^2 \cdot C}= 1.20~\mathrm{H}.\]


Now, knowing the value of $f_{анти} = 1.999740~\mathrm{MHz}$, we can get the value of $L$:


\[L = \dfrac{1}{4 \pi^2 f_\text{анти}^2 \cdot C}= 1.20~\mathrm{H}.\]

\par\smallskip\noindent\textbf{Answer:} \textbackslash{}textbf{Answer:} \[L = 1.20~\mathrm{H}; \]\[C = 5.3 ~фФ;\]\[C_0 = 53.9~\mathrm{pF}.\]
\end{solcard}

\begin{solcard}{A12}{From the results of points A4, A6, A10, get the value $R$.}{1.20}

Let us first obtain the answer with the approximation $\dfrac{4 \omega_p^2 \alpha^2}{(\omega_p^2-\omega_s^2)^2} \ll 1,$ which follows from the condition by simply expressing $\alpha$ from the formula obtained in paragraph A8. In this case $\alpha = 164.4~\dfrac{\text{\Omega}}{\text{H}}$.


Note. It is possible to find $R$ without approximation from the condition (this was not required from the participants). In paragraph A6, we received that $\alpha = \dfrac{R}{L},$ remains to be determined $\alpha.$. The easiest way to determine it is by resonance $Z.$ As has already been found: \[\omega = \sqrt{\dfrac{(\omega_s^2+\omega_p^2)\pm \sqrt{(\omega_s^2+\omega_p^2)^2+4\omega_p^2 (\alpha^2 - \omega_s^2)} }{2}}\] for extrema $|Z|$. Let's substitute them into the expression for $|Z|^2:$ $$|Z|_\text{крит}^2 = \dfrac{\left(\omega_s^2 - \dfrac{(\omega_s^2+\omega_p^2)\pm \sqrt{(\omega_s^2+\omega_p^2)^2+4\omega_p^2 (\alpha^2 - \omega_s^2)} }{2}\right)^2 + \left(\omega_p \alpha \right)^2}{(\omega_p C_0)^2 \left( \left(\omega_p^2 - \dfrac{(\omega_s^2+\omega_p^2)\pm \sqrt{(\omega_s^2+\omega_p^2)^2+4\omega_p^2 (\alpha^2 - \omega_s^2)} }{2} \right)^2 + \left(\omega_p \alpha \right)^2 \right)},$$ here signу $+$ соответствует $|Z|_\text{max}^2,$ signу $-$ соответствует $|Z|_\text{min}^2.$ Transforming данное expression: $$|Z|_\text{крит}^2 = \dfrac{ \left( \omega_p^2-\omega_s^2 \pm \sqrt{(\omega_p^2-\omega_s^2)^2+4 \omega_p^2 \alpha^2)}\right)^2+4 \omega_p^2 \alpha^2}{\left( \omega_p^2-\omega_s^2 \mp \sqrt{(\omega_p^2-\omega_s^2)^2+4 \omega_p^2 \alpha^2)}\right)^2+4 \omega_p^2 \alpha^2} \cdot \dfrac{1}{\omega_p^2 C_0^2},$$ $$|Z|_\text{крит}^2 = \dfrac{ \left( 1 \pm \sqrt{1+\dfrac{4 \omega_p^2 \alpha^2}{(\omega_p^2-\omega_s^2)^2})}\right)^2+\dfrac{4 \omega_p^2 \alpha^2}{(\omega_p^2-\omega_s^2)^2}}{\left( 1 \mp \sqrt{1+\dfrac{4 \omega_p^2 \alpha^2}{(\omega_p^2-\omega_s^2)^2})}\right)^2+\dfrac{4 \omega_p^2 \alpha^2}{(\omega_p^2-\omega_s^2)^2}} \cdot \dfrac{1}{\omega_p^2 C_0^2}.$$ Введём обоvalue $y = \dfrac{4 \omega_p^2 \alpha^2}{(\omega_p^2-\omega_s^2)^2}:$ $$|Z|_\text{крит}^2 = \dfrac{ \left( 1 \pm \sqrt{1+y}\right)^2+y}{\left( 1 \mp \sqrt{1+y}\right)^2+y} \cdot \dfrac{1}{\omega_p^2 C_0^2}.$$ Оlimitяем $|Z|_{\max}$ мы госдо точнее, чем $|Z|_{\min}$, therefore по $|Z|_\text{max} \approx 12500~\text{Ohm},$ найденное по графику в partе A4 let us find $y$: $$|Z|_\text{max}^2 \cdot \omega_p^2 \cdot C_0^2 = 71.67 = \dfrac{ \left( \sqrt{1+y}+1\right)^2+y}{\left( \sqrt{1+y}-1\right)^2+y}.$$ Решим данное equation for/at помощи численных методов (метод итераций, таблицы в калькуляторе or любой другой рабочий метод): $$y = 0.05740 \Rightarrow \alpha = 165.6~\dfrac{\mathrm{\Omega}}{\mathrm{H}} \Rightarrow R  = L \cdot \alpha = 199~\mathrm{\Omega},$$ which is quite close to the approximate answer obtained above, that is, the approximation from the condition for a given quartz resonator is correct.

\par\smallskip\noindent\textbf{Answer:} \textbackslash{}textbf{Answer:} \[R = 197~\mathrm{\Omega}.\]
\end{solcard}

\end{document}